\scanpage{37}
\begin{center}
UNIVERSITÉ DE PARIS

\figure{\begin{tikzpicture}\draw (-.4,0)--(-.06,0);\fill (0,0) circle (.7pt);\draw (.06,0)--(.4,0);\end{tikzpicture}}

ÉCOLE NORMALE SUPÉRIEURE

\figure{\begin{tikzpicture}\draw (-.4,0)--(-.06,0);\fill (0,0) circle (.7pt);\draw (.06,0)--(.4,0);\end{tikzpicture}}

45, RUE D'ULM

ODÉON 06$\cdot$45

\figure{\begin{tikzpicture}\draw (-.4,0)--(-.06,0);\fill (0,0) circle (.7pt);\draw (.06,0)--(.4,0);\end{tikzpicture}}
\end{center}

Il résulte du lemme 1 que $\chi$ est surjective.

---

Montrons que $\chi$ est \deleted{\unsure{aussi}} un homomorphisme : \inserted{Soient $g_1,g_2\in G$} Soit $\gamma_1$ un \deleted{\unsure{la}} chemin tel que $\partial_1(\gamma_1)=x_0$, $\partial_0(\gamma_1)=g_1x_0$ ; idem p\textsuperscript{o} $\gamma_2$ tel que $\partial_1(\gamma_2)=x_0$, $\partial_0(\gamma_2)=g_2x_0$. On a
\[
\partial_1(g_1\gamma_2)=g_1x_0\text{ et }\partial_0(g_1\gamma_2)=g_1g_2x_0:
\]
le chemin \inserted{$\gamma=$}$(\gamma_1,g_1.\gamma_2)$ vérifie $\partial_1\gamma=x_0$, $\partial_0\gamma=g_1g_2x_0$, \deleted{\unsure{on a}} On a
\[
p\gamma=(p\gamma_1,pg_1\gamma_2)=(p\gamma_1,p\gamma_2)\quad\deleted{\unsure{d'où}}\quad\text{d'où}\quad\chi(g_1g_2)=\chi(g_1)\chi(g_2)
\]

---

Le noyau de $\chi$ contient évidemment le stabilisateur de $x_0$. Soit $x$ un sommet de $E$ et $s\in G$ tel que $sx=x$. Soit $\lambda$ un chemin de $E$ tel que $\partial_1(\lambda)=x_0$, $\partial_0\lambda=x$.
\deleted{\unsure{le lacet}} $\gamma=(\lambda,-s\lambda)$ est un chemin de $E$ tel que $\partial_1(\gamma)=x_0$, $\partial_0\gamma=sx_0$.
d'où
\[
\chi(s)=\text{classe }(p(\lambda,-s\lambda))=\text{classe }(p(\lambda,-\lambda))=e\text{ (élément neutre)}.
\]
Le noyau de $\chi$ contient donc le sous groupe $S$ \deleted{\unsure{distingué}} engendré par les stabilisateurs des \deleted{\unsure{points}} sommets de $E$. Ce sous groupe est distingué puisque le conjugué du stabilisateur de $x$ est le stabilisateur d'un transformé de $x$.

---

On sait que $\operatorname{Ker}\chi\supset S$. Montrons qu'il ne le contient pas strictement.

\uline{Lemme 3} : $E/S$ est un revêtement\inserted{normal} de $E/G$, de groupe $G/S$.
\mspage{36}

